\subsection{Sensitivity to the calibration rule}
\label{subsec:calibration_sensitivity}

To examine whether the calibration depends critically on a single
shrinkage setting, we considered
\[
\theta_{\kappa}=
\begin{cases}
(\Delta-\kappa s)_{+}/\Delta,&\Delta>0,\\
0,&\Delta\le0,
\end{cases}
\]
with \(\kappa\in\{0,0.5,1,1.5\}\), and varied the number of
stratified folds over \(F\in\{3,5,10\}\). The value \(\kappa=0\)
reduces to hard graph selection, whereas the proposed method uses
\(\kappa=1\). PIE, YaleB, and COIL100 were fixed before this experiment
to represent fallback, borderline, and clear-benefit regimes,
respectively. Every configuration used the same labeled subsets and
initial factors within each paired trial.

% Insert:
% CGC_GOCNMF_sensitivity3_20seed_kappa_table.tex
% CGC_GOCNMF_sensitivity3_20seed_fold_table.tex

\subsection{Six-dataset ablation}
\label{subsec:ablation6}

We compared four graph variants on all six benchmark datasets:
BASE uses \(W^{(0)}\), LOCAL uses \(W^{(1)}\), HARD selects
\(W^{(1)}\) when \(\Delta>0\) and otherwise uses \(W^{(0)}\), and
SHRINK uses the proposed calibrated graph. All variants shared the same
labeled subsets, factor initialization, graph support, regularization
coefficient, and iteration count.

% Insert:
% CGC_GOCNMF_ablation6_20seed_table.tex
% CGC_GOCNMF_ablation6_20seed_calibration_table.tex
